\phantomsection
\addcontentsline{toc}{chapter}{Introduction}
\begin{center}
    {\scshape\bfseries\Large Introduction}
\end{center}
\markboth{INTRODUCTION}{}
\thispagestyle{introduction}
In 1965 M.S. Narasimhan and C.S. Seshadri proved the so-called Narasimhan-Seshadri theorem \parencite{NS1965}. In their original formulation, their result provided a relation between the notion of stability for holomorphic vector bundles over a compact Riemann surface, which had been introduced by D. Mumford in 1963 in the context of Geometric Invariant Theory, and representations of the fundamental group of the surface. Briefly, stability of vector bundles is defined by assigning a numerical invariant to each bundle, called its slope. A vector bundle over a Riemann surface is stable if its slope is greater than the slope of any of its proper holomorphic subbundles. In the 1970s and 1980s, the works by M. Atiyah, R. Bott, S. Kobayashi, S. Donaldson, N. Hitchin and others showed that the result by Narasimhan and Seshadri could be translated to the language of differential geometry, and more specifically, in terms of Yang-Mills connections, notions arising from aspects of differential geometry and theoretical physics that were being studied at the time. More precisely, Donadlson gave a new proof of the theorem by Narasimhan and Seshadri in 1983 \parencite{DonaldsonNS83}. He proved that an indecomposable holomorphic vector bundle over a compact Riemann surface \(\mathcal{E} \to X\) is stable if and only if its hermitian fiber metric \(h\), whose Chern connection \(\nabla\) satisfies the equation \(\star F_\nabla=-\ii c \cdot \mathrm{id}_{\mathcal{E}}\) for a fixed constant \(c\) that depends on the topology of the bundle and the surface. This is called the \textit{constant central curvature} condition. If the condition of being decomposable is omitted, one gets that a hermitian holomorphic vector bundle is polystable if and only if it admits a hermitian fiber metric whose Chern connection satisfies the constant central curvature condition. The relation with Yang-Mills theory can be seen even more directly if one realizes that the theorem by Donaldson is equivalent to the following statement: a smooth complex vector bundle \(E \to X\) over a compact Riemann surface endowed with a hermitian fiber metric has a \textit{Hermite-Einstein-Yang-Mills connection}, i.e. a connection satisfying \(\star F_\nabla=-\ii c \cdot \mathrm{id}_E\) if and only if it admits a holomorphic structure making the bundle polystable.
\\

In 1980, Kobayashi generalized the notion of stability to holomorphic vector bundles over higher-dimensional compact Kähler manifolds. This was a remarkable contribution, as it required the introduction of coherent sheaves. Kobayashi and Hitchin conjectured independently that the theorem of Narasimhan and Seshadri holds for compact Kähler surfaces (hence the name, \textit{Hitchin-Kobayashi correspondence}). That is, a holomorphic vector bundle over a compact Kähler manifold is polystable if and only if it admits a hermitian fiber metric whose Chern connection satisfies the equation \(\ii \Lambda_\omega F_\nabla=c\cdot \mathrm{id}_E\) where \(\Lambda_\omega\) is the dual Lefschetz operator. Such a fiber metric is called a Hermite-Einstein metric. The fact that the existence of the metric implies polystability is relatively easy to prove. The converse, however, was remarkably difficult, and it was proved by Donaldson and K. Uhlenbeck and S.T. Yau in successive works in the 1980s. This result is surprising because it provides a direct relation between a purely algebraic notion, stability, and a differential-geometric notion, the existence of a Hermite-Einstein metric. Again, this can be stated equivalently by saying that a complex vector bundle \((E,h)\to X\) has a unitary connection satisfying the \textit{Hermite-Yang-Mills equations}
\[
\left\{
\begin{array}{ccc}
    F_\nabla^{(0,2)} & = & 0 \\
    \Lambda_\omega F_\nabla & = & -\ii c \cdot \mathrm{id}_E
\end{array}
\right.
\]
if and only it has a holomorphic structure making it polystable. \\

This correspondence acts like a dictionary in the sense that it allows us to translate algebraic problems into differential-geometric problems and vice versa. The theorem above can be interpreted as fixing a hermitian fiber metric on a complex smooth vector bundle. Then, it can be interpreted in terms of moduli spaces: there is a homeomorphism between the moduli spaces of polystable holomorphic structures on the vector bundle and the moduli space of \textit{Hermite-Yang-Mills connections} on \(E\). We insist on the fact that this result is extraordinary, in the sense that it is a purely gauge-theoretic approach to an algebraic notion.\\

In 1975, the result by Narasimhan and Seshadri was generalized by Ramanathan for principal bundles over a compact Riemann surface with compact connected reductive complex Lie groups \parencite{ramanathan1975stable}. Among the contributions of Ramanathan to the discipline, we should highlight the fact that he gave a notion of stability. This was generalized for compact Kähler manifolds (which again required modifying the notion of stability) by Ramanathan and Subramanian \parencite{SubramanianRamanathan}. From the beginning of the century onwards, several generalizations of the correspondence have been proven. Namely, I. Mundet i Riera gave a generalization to Kähler fibrations \parencite{i1999hitchin}, and O. García Prada and P.B. Gothen gave a generalization of the correspondence for Higgs bundles \parencite{garcia2009hitchin}. \\

The remaining piece of the puzzle was given by M. Atiyah and R. Bott, who published a ground-breaking paper \parencite{atiyah1983yang} in 1983 that set the bases of the moment map picture. Basically, they showed that the moduli space of Hermite-Einstein connections could be described by means of a symplectic reduction for an appropriate moment map on the set of connections, which is acted on by the gauge group of the bundle. Phrased that way, the Narasimahn-Seshadri theorem and, more generally, the Hitchin-Kobayashi correspondence, can be thought of as infinite-dimensional versions of the Kempf-Ness theorem. %The Kempf-Ness theorem is treated in chapter \ref{chapter_kempf_Ness}. It relates symplectic reduction and GIT quotients, leading to a shocking relation between a purely differential construction: symplectic reduction, and a purely algebraic one: geometric invariant theory. An infinite-version of this result is just as surprising. \\

The aim of this thesis is to present an account of the Hitchin-Kobayashi correspondence, explaining the moment map picture and the sense in which it can be seen as an infinite-dimensional version of the Kempf-Ness theorem, as well as the way in which the finite-dimensional version of the Kempf-Ness theorem motivates the correspondence. Detailed proofs of these results are beyond the scope of this thesis. However, for completeness, we will try to sketch them. We will first give a brief account of the vector bundle version of the correspondence both, to motivate the principal bundle, and so that we can compare it with the principal bundle version. We will, however, devote our efforts mostly to the principal bundle version of the correspondence. We have followed an exposition close to that of \parencite{i1999hitchin} and \parencite{garcia2009hitchin}, but making the corresponding simplifications (which basically consist of ignoring any extra structure). When studying moduli spaces for these structures, rigorous constructions require working with Sobolev completions of several spaces, something that is often dismissed in the literature omitting the technical details. We have tried to specify, at least for the principal bundle version of the theorem, why Sobolev completions are needed to make everything rigorous and how to work with them. The second objetive of the thesis is to set the first steps towards a generalization of the correspondence for principal bundles over non-connected structure groups. Again, this requires both a new Hermite-Einstein-Yang-Mills equation and a new notion of stability, but more importantly, it relies on the work by García Prada, G. Barajas, Gothen and Mundet i Riera in \parencite{barajas2023non} on twisted equivariant bundles over coverings. The key ideas and results concerning this generalization have been taken from a draft of a yet unpublished article that my supervisor kindly shared with me. \\

\thispagestyle{introduction}
As a brief outline of the thesis, in chapter \ref{chapter_kempf_Ness} we address the finite-dimensional version of symplectic reduction and the Kempf-Ness theorem. The former is generalized to Banach manifolds in chapter \ref{Chapter_inf_dim_sympl_red}. Chapters \ref{ch_connections} and \ref{chapter_principal_bundles} introduce connections in the context of vector and principal bundle, and they treat the relation between connections and complex structures; they so-called Chern or Chern-Singer correspondences. Finally, chapter \ref{chapter_HK_correspondence} builds on all the previous work and introduces the Hitchin-Kobayashi correspondence.\\

I have tried to make this thesis as self-contained as possible, an ultimately impossible task, given the sheer amount on concepts that play a role in the Hitchin-Kobayashi correspondence. For the convenience of the reader, and for my own convenience when citing results, I have included 3 appendices that give short introductions to some of the prerequisites needed for the main body of the text. Appendix \ref{appendix_Kahler} is a short introduction to the theory of complex manifolds and vector bundles. It should be read before the main body of the text. Appendix \ref{appendix_lie} contains some notions of the theory of Lie groups which may not be covered in a first course on the subject, as well as some tools from \parencite{barajas2023non} which set the basic framework needed to work with principal bundles with non-connected structure group. It should be read before chapter \ref{chapter_HK_correspondence} although some of its results are used earlier in the text. Finally, appendix \ref{appendix_Banach_manifolds} gives a short introduction to Banach manifolds. It mostly contains relevant definitions, and statements that show that Banach manifolds are sufficiently well-behaved to still satisfy much of the properties of usual finite-dimensional smooth manifolds (with the exception, of course, of the existence of charts of the form \((x^1, \ldots x^n)\)). It should be read before chapter \ref{Chapter_inf_dim_sympl_red}. With my deepest sorrow, and in an effort to keep the length of the thesis manageable, I have omitted a lot of proofs and only sketched others. However, whenever a proof has been omitted, I have tried to give a reference in the literature where it can be found. On the other hand, I have tried to pay special attention to detail in the proofs that I have included. This is a matter of personal taste, but I hope it will be helpful to future readers. I have also included an alphabetical index and an index of notations at the end for the convenience of the reader. \\

%\addcontentsline{toc}{chapter}{Introduction}